\section{A degree--amplitude inequality from random restrictions}
\label{sec:random-restriction}
The next lemma uses the global margin to strengthen a degree-only
obstruction. Its proof combines random restrictions with the parity
approximation bound of Aspnes, Beigel, Furst, and Rudich.

\begin{lemma}[Robust degree--amplitude inequality]
\label{cor:restriction-amplitude}
\label{lem:random-restriction-exact}
Let $f=\chi_U$ on the uniform $d$-cube, where $|U|=k$, and let
$|z(x)|\ge a>0$ everywhere. Suppose $z=P+T$, where $\deg P\le u$,
$1\le u\le k$, and $P,T$ have disjoint Fourier supports. There are
absolute constants $c,C>0$ such that
\begin{equation}
 \Pr\{\operatorname{sign}z\ne f\}
 \ge \frac7{16}-\frac{\|T\|_2^2}{a^2}
       -C\frac{\|P\|_2^2}{a^2}\exp\!\left(-c\frac{k}{u}\right).
 \label{eq:restriction-amplitude}
\end{equation}
More precisely, for integers $1\le m\le k$ and $0\le q<m$ the right-hand
side can be replaced by
\begin{equation}
 2^{-m}\sum_{j=0}^{\lfloor(m-q-1)/2\rfloor}\binom mj
 -\frac{\|T\|_2^2}{a^2}
 -\frac{\|P\|_2^2}{a^2}
    \min\!\left\{1,
       \left(\frac{eum}{k(q+1)}\right)^{q+1}\right\}.
 \label{eq:restriction-exact}
\end{equation}
\end{lemma}
\begin{proof}
Choose a uniform $m$-element subset $J\subseteq U$, fix the coordinates
outside $J$ uniformly, and truncate the restricted polynomial $P$ to
degree $q$ in the remaining coordinates. Denote the resulting polynomial
by $Q$. Viewed on the original cube, this operation retains exactly the
Fourier terms $\widehat P(S)\chi_S$ with $|S\cap J|\le q$. Orthogonality
therefore gives
\[
 \mathbb E_J\|z-Q\|_2^2
 =\|T\|_2^2+
   \sum_S\widehat P(S)^2\Pr_J\{|S\cap J|>q\}.
\]
Since $|S|\le u$, the probability is bounded by the upper tail of a
hypergeometric random variable with population size $k$, at most $u$
marked elements, and sample size $m$. The factorial-moment estimate gives
\[
 \Pr\{|S\cap J|\ge q+1\}
 \le \min\!\left\{1,
       \left(\frac{eum}{k(q+1)}\right)^{q+1}\right\}.
\]
For every restriction, the target is parity or its negation on the
$m$ remaining coordinates. Lemma~\ref{lem:ptf-parity} gives error
at least
\[
 B(m,q):=2^{-m}\sum_{j=0}^{\lfloor(m-q-1)/2\rfloor}\binom mj
\]
for every degree-$q$ polynomial threshold function on that cube. This
statement holds for real coefficients: its proof uses real linear
algebra. Ties can be assigned either fixed sign by adding an arbitrarily
small constant. Whenever $z$ and $Q$ have different signs, the global
margin of $z$ implies $|z-Q|\ge a$. Averaging the resulting error bound
proves~\eqref{eq:restriction-exact}.

For completeness, choose a sufficiently small absolute $\delta>0$,
put $t=k/u$, and, when $t$ exceeds a sufficiently large absolute constant,
set $m=\min\{k,\lfloor\delta t^2\rfloor\}$. If $m=k$, then
$u\le\sqrt{\delta k}$; use $Q=P$ directly. The central binomial
atom bound yields $B(k,u)\ge7/16$ after increasing the fixed lower
threshold on $t$. If $m<k$, take $q=\lfloor\sqrt m/32\rfloor$.
For example $\delta=(1024e)^{-2}$ ensures
\[
 \frac{eum}{k(q+1)}\le32e\sqrt\delta\le\frac1{32},
 \qquad q+1\ge c_0 t,
 \qquad B(m,q)\ge\frac7{16}
\]
for sufficiently large $t$, with $c_0>0$ absolute. This proves
\eqref{eq:restriction-amplitude} in that range. For bounded $t$,
use $\|P\|_2^2+\|T\|_2^2=\|z\|_2^2\ge a^2$ and enlarge $C$ so
that the asserted lower bound is nonpositive.
\end{proof}

In particular, any degree-$u$ polynomial with global margin $a$ whose
threshold has parity error at most $3/8$ must satisfy
\[
 \|P\|_2^2\ge c' a^2\exp(c k/u).
\]
For a fixed degree fraction $u/k$, this lower bound is a constant. It is
therefore consistent with the bounded-amplitude blocked construction in
Proposition~\ref{prop:blocked-approximants}.

% Primary reference to add to the bibliography:
% James Aspnes, Richard Beigel, Merrick Furst, and Steven Rudich.
% The expressive power of voting polynomials. Combinatorica 14(2), 1994.
% Theorem 2.2 and Lemma 2.3.
% https://www.cs.yale.edu/homes/aspnes/papers/stoc91voting.pdf
